\subsection{The actual central word realizing the quadratic coefficient}
\label{sec:nonstandard-central-word}

\paragraph{P-F52 (an intertwiner inside the displayed nonstandard algebra).}
Use the cubic algebra $B_0$ over $F_0=\mathbb C(q)$, its splitting
field $K=F_0(x)$, and the elements $c,c',P,D,F$ of P-F36 and P-F39.
Here $x^2=D$; the letter $F$ without a subscript continues to denote
the positive Laurent normalization in (P-F39a).
Write $A=Q_1$, $B=Q_2$, and
\[
 f=(q^4+1)(q^4-q^2+1)(q^2+1)/q^5
   =q^5+q+q^{-1}+q^{-5}.
\]
Define the following explicit Laurent polynomial and central word:
\[
 \begin{split}
 S={}&q^{10}+q^{-10}+3(q^8+q^{-8})+(q^6+q^{-6})
       +5(q^4+q^{-4})-2(q^2+q^{-2})+16,\\
 Z={}&(Q_1-Q_2)^2+(S-f^2)1.
 \end{split}                                                    \tag{P-F52a}
\]
In the actual source blocks $U_1,U_2\subset B_0\otimes K$,
left and right multiplication by $Z$ have scalar values $-c,-c'$,
respectively. In particular, this operator is specified by a short
word in the original generators, rather than by assigning arbitrary
eigenvalues to abstract central idempotents.

There is an explicit two-dimensional $F_0$-subspace of $B_0$ on
which multiplication by $Z$ has precisely the matrix $M_0$ of
(P-F39k). Its multiple $FZ$ has matrix $M=FM_0$.
The same construction gives four copies of this matrix on the
eight-dimensional descended two-sided ideal associated with
$U_1\oplus U_2$. These are statements about actual multiplication
in $B_0$, with the original source basis retained.

For every real $q>0$, the same word $Z$ acts as a strictly positive,
self-adjoint operator in the source's actual tensor representation
on $X_q^{\otimes3}$, with the explicit positive tensor form given
below. So does $FZ$. This last assertion concerns the tensor
specialization, whose image can have smaller dimension at exceptional
parameters; the generic $F_0$-linear intertwiners are not asserted to
remain injective at every specialization.

\paragraph{Only the original relation and the original cell element are needed.}
The source gives $A^2=fA$, $B^2=fB$, and the independent
alternating-word basis of lengths at most ten: GCT VII,
Section~7.1.1, pp.~18--19. Put
\[
       C=(A-B)^2=f(A+B)-AB-BA,
       \qquad G=f^2-C,
       \qquad w_j=A(BA)^j.
                                                               \tag{P-F52b}
\]
Direct multiplication gives
$AC=CA=f^2A-ABA$ and $BC=CB=f^2B-BAB$.
Thus $C$ and $G$ are central, and
\[
                    Gw_j=w_{j+1}\qquad(j\ge0).                 \tag{P-F52c}
\]
The actual relation in Figure~2, printed p.~31, has the form
\[
            w_5-Sw_4+a_3w_3+a_2w_2+a_1w_1+a_0w_0=0,          \tag{P-F52d}
\]
where the coefficient of $w_5$ is one and the coefficient of
$w_4$ is exactly $-S$ from (P-F52a). The actual element $u=u_1^1$ in
Figure~7, printed p.~36, has the form
\[
                  u=w_4+cw_3+b_2w_2+b_1w_1+b_0w_0.            \tag{P-F52e}
\]
Its conjugate $u'=u_2^1$ replaces $x$ by $-x$ and therefore has
coefficient $c'$ at $w_3$. This is the same scalar coefficient
used in GCT VIII and P-F39, even though VIII also substitutes
$P$-words for $Q$-words in its corresponding element.

The following full coefficient identity verifies the eigenvector
equation without relying on the simple-block assertion. Set
\[
 \begin{gathered}
 R=-3q^{28}-4q^{24}+2q^{22}-10q^{20}+2q^{18}-4q^{16}-3q^{12},\\
 g_1=\frac{-R-x}{2q^{20}},\qquad
 g_2=\frac{-R+x}{2q^{20}},\\
 g_3=q^4+2+q^{-4},\quad g_4=q^6+2+q^{-6},\quad
 g_5=q^{10}+2+q^{-10},\\
 E_1=g_3+g_4+g_5,\quad E_2=g_3g_4+g_3g_5+g_4g_5,\quad
 E_3=g_3g_4g_5,\\
 L=-R/q^{20},\qquad
 N_0=\frac{(q^4-q^2+1)^2
             (q^{12}+q^{10}+2q^8+2q^4+q^2+1)^2}{q^{16}}.
 \end{gathered}                                                 \tag{P-F52f1}
\]
Expansion gives $R^2-D=4q^{40}N_0$, $g_1+g_2=L$,
$g_1g_2=N_0$, and $L+E_1=S$. All coefficients of the
original relation and the original cell element can be written
as follows; the expressions are direct factorizations of
Figures~2 and 7, including their lower coefficients:
\[
\begin{array}{c|c|c}
 j&[t^j]r(t)&[t^j]p_1(t)\\ \hline
 0&-N_0E_3&g_2E_3\\
 1&N_0E_2+LE_3&-E_3-g_2E_2\\
 2&-N_0E_1-LE_2-E_3&E_2+g_2E_1\\
 3&N_0+LE_1+E_2&-E_1-g_2=c\\
 4&-L-E_1=-S&1\\
 5&1&0
\end{array}                                                     \tag{P-F52f2}
\]
Here $r(G)A=0$ is exactly (P-F52d), and $p_1(G)A=u$
is exactly (P-F52e). For example their constant coefficients
are the factored first rows of the two source figures; the
other rows follow by the same multiplication and collection
of powers of $q$ and $x$. The accompanying exact verifier
transcribes and compares all six relation coefficients and
all five cell coefficients, without interpolation.
Multiplying the polynomials in the table now gives
\[
 \begin{split}
 r(t)&=(t^2-Lt+N_0)(t^3-E_1t^2+E_2t-E_3),\\
 p_1(t)&=(t-g_2)(t^3-E_1t^2+E_2t-E_3),\\
                 (t-g_1)p_1(t)&=r(t).
 \end{split}                                                    \tag{P-F52f3}
\]
Substitution of the central element $G$ proves
$(G-g_1)u=r(G)A=0$ directly. Since $g_1=S+c$,
\[
                 Gu=(S+c)u,\qquad Zu=-cu.                     \tag{P-F52f}
\]
The leading coefficient gives an independent quick check:
in $Gu$ only the reduction of $w_5$ and the term $cw_4$
contribute to $w_4$, with coefficient $S+c$.
Conjugating gives $Zu'=-c'u'$. Both elements are nonzero because
their coefficient at the independent word $w_4$ is one.
Moreover $c-c'=-x/q^{20}\ne0$, so they are distinct; their
belonging to different direct summands also makes them
linearly independent over $K$.

To label the corresponding full blocks, use the split Wedderburn
decomposition in GCT VII, Section~7.1.2, pp.~19--20: $u,u'$
belong to the four-dimensional simple matrix-algebra summands
$U_1,U_2$. The center of a matrix algebra consists of scalars,
as follows by commuting with its diagonal and off-diagonal matrix
units. Thus (P-F52f) gives the same scalar on the full respective
summands, on either side. The eigenvector identity itself has
already been verified independently in (P-F52f1)--(P-F52f3).

This argument does not use the printed two-by-two generator
matrices on p.~20. As printed, both are upper triangular and
would give a common invariant line, inconsistent with the
generic irreducibility claimed there. That display is
quarantined here; no silent correction to it is needed for
(P-F52f), which follows from the complete original relation
and cell-element factorization. The cited split simple-algebra
assertions serve separately to identify the full block labels.

\paragraph{The exact descended intertwiner.}
Let $\tau(x)=-x$ be the nontrivial automorphism of $K/F_0$.
Set
\[
            v_0=u+u',\qquad v_1=x(u-u').                       \tag{P-F52g}
\]
They are fixed by $\tau$ and hence lie in $B_0$. To see this
last assertion explicitly, expand in the $F_0$-word basis:
each fixed coefficient in $K$ belongs to $F_0$.
The coefficient of $w_4$ in $v_0$ is two and in $v_1$ is zero,
while that of $w_3$ in $v_1$ is
$x(c-c')=-D/q^{20}\ne0$. Thus $v_0,v_1$ are independent.

The $F_0$-linear map
\[
       \iota:K\longrightarrow B_0,
           \qquad a\longmapsto au+\tau(a)u'                  \tag{P-F52h}
\]
is injective and sends the basis $(1,x)$ to $(v_0,v_1)$.
Equation (P-F52f) gives
\[
       Z\,\iota(a)=\iota\left(\frac{P+x}{2q^{20}}a\right).
                                                               \tag{P-F52i}
\]
In that basis, actual multiplication by $Z$ is therefore
\[
                 \frac1{2q^{20}}
                       \begin{pmatrix}P&D\\1&P\end{pmatrix}, \tag{P-F52j}
\]
exactly (P-F39k), with no transpose, sign, or scalar discrepancy.
Multiplying $Z$ by the source normalization $F$ gives (P-F39m).
The map $\iota$ intertwines the indicated central action; its
two-dimensional image is not asserted to be a left ideal of $B_0$.

For the actual two-sided ideal, take the four source basis elements
of $U_1$, namely $u_1^1,u_1^{12},u_1^{21},u_1^2$, and their
conjugates in $U_2$. Call them $u_r,u'_r$, respectively.
The eight vectors $u_r+u'_r$, $x(u_r-u'_r)$ are $F_0$-rational
and independent: a relation, projected onto $U_1$ and $U_2$,
has coefficients $a_r+xb_r=0$ and $a_r-xb_r=0$, forcing
$a_r=b_r=0$. Conversely every fixed vector in
$U_1\oplus U_2$ is of this form by writing its $U_1$ coefficients
as $a_r+xb_r$. Their span is therefore the descended two-sided
ideal $W=(U_1\oplus U_2)^\tau$, of dimension eight over $F_0$.
Multiplication by $Z$ acts on it as four copies of (P-F52j).
The $U_i$ here are matrix-algebra summands, of dimension four;
the source's simple left modules $T_i$ have dimension two.
These two roles are not identified.

\paragraph{An explicit positive form for the actual tensor action.}
For $q>0$ set
\[
 [2]=q+q^{-1},\quad [3]=q^2+1+q^{-2},\qquad
 (d_0,d_1,d_2,d_3)
       =\bigl(1,q^{-1}[3],[3][2]^2,q^3[3]^2[2]^2\bigr).
                                                               \tag{P-F52k}
\]
All four entries are positive. Give the source basis
$x_i=f^ix_0$ of $X_q$ the diagonal Gram matrix
$\operatorname{diag}(d_0,d_1,d_2,d_3)$, and use its tensor
product on tensor powers. Here the lowering operator denoted
$f$ in the source's basis notation is distinct from the scalar
Laurent polynomial $f$ in (P-F52a).

The actual rescaled projection matrix $\mathsf P$ in GCT VII,
Figure~1, printed p.~30, is self-adjoint for this tensor form.
The entire check can be organized by $i+j$, since its nonzero
entries preserve that sum. The one-dimensional blocks and the
two-dimensional blocks $i+j=1,5$ have real diagonal entries
and equal opposite off-diagonal entries, with equal Gram weights.
For $i+j=2$, the weights are $(d_2,d_1^2,d_2)$; the only
nontrivial opposite-entry ratio is verified by
\[
 \frac{d_2}{d_1^2}
       =\frac{q^2(q^2+1)^2}{q^4+q^2+1}.
\]
For $i+j=3$, the weights are
$(d_3,d_1d_2,d_1d_2,d_3)$; the displayed outer-to-inner
opposite-entry ratios are $q^{-4}$, and
$d_3/(d_1d_2)=q^4$. Opposite entries between equal-weight
positions are equal. Finally, for $i+j=4$ the weights are
$(d_1d_3,d_2^2,d_1d_3)$, and the required ratio is
\[
 \frac{d_1d_3}{d_2^2}
       =\frac{q^2(q^4+q^2+1)}{(q^2+1)^2}.
\]
Substitution in the entries of Figure~1 gives
$d_id_j\mathsf P_{ij,kl}=d_kd_l\mathsf P_{kl,ij}$
in every block. These identities are Laurent identities,
so also hold at $q=1$, where some off-diagonal entries vanish;
no division by such an entry is being asserted there.
Thus $\mathsf P$ is self-adjoint. The source's complementary
rescaled projection is $\mathsf Q=fI-\mathsf P$ by
equations~(4), (16)--(18), pp.~6--7 and 18, so it too is
self-adjoint. Consequently $Q_1=\mathsf Q\otimes I$ and
$Q_2=I\otimes\mathsf Q$ on $X_q^{\otimes3}$ are
self-adjoint for the displayed positive product form.

Put $t=q^2+q^{-2}=2+v$, so $v\ge0$. A direct expansion gives
\[
 \begin{split}
 S-f^2
 &=3(q^8+q^{-8})-(q^6+q^{-6})
       +3(q^4+q^{-4})-3(q^2+q^{-2})+12\\
 &=3t^4-t^3-9t^2+12
   =3v^4+23v^3+57v^2+48v+16>0.
 \end{split}                                                    \tag{P-F52l}
\]
For every nonzero tensor vector $w$, self-adjointness now gives
\[
 \langle w,Zw\rangle
   =\|(Q_1-Q_2)w\|^2+(S-f^2)\|w\|^2>0.
\]
This proves actual positive-definiteness, without importing a
generic compact-star assertion across an exceptional specialization.
The coefficients of $F$ are positive, so $F(q)>0$ and the same
conclusion holds for $FZ$.\hfill$\square$

\paragraph{What this finite comparison settles.}
The positive quadratic operator of P-F39 is now realized by an
explicit element of the actual nonstandard algebra, with a
source-basis intertwiner and an explicit tensor positive form.
In particular the occurrence of the eigenvalue $-c$ is tied to
the displayed word coefficient of $u_1^1$ through the leading-term
calculation (P-F52d)--(P-F52f); it is not imposed after choosing an
unrelated matrix model. On the descended two-dimensional space,
however, the unscaled determinant still has coefficients $T,N$
from (P-F39c), with their negative Laurent coefficients. Thus the
actual positivity just proved does not imply coefficientwise
positivity. The prescribed scale $F$, the other source coefficients,
their canonical-basis meanings, and any cohomological or theta
intertwiner remain substantive additional data. No uniform positivity
theorem in tensor degree or PH2/SH consequence is claimed here.

Primary source: GCT VII, Section~7.1.1--7.1.2, pp.~18--20,
Figure~1 p.~30, Figure~2 p.~31, Figure~7 p.~36,
\url{https://arxiv.org/pdf/0709.0749}; GCT VIII's matching scalar
coefficient and its normalization are cited and checked in P-F39,
\url{https://arxiv.org/pdf/0709.0751}.

\paragraph{A finite normalization corollary to P-F52.}
Even the combination of positive actual central action and symmetric
Tate--Lefschetz trace and norm does not characterize the source's
prescribed normalization $F$. Indeed put
\[
       H(q)=(q+1+q^{-1})^4
        =19+16(q+q^{-1})+10(q^2+q^{-2})
                    +4(q^3+q^{-3})+(q^4+q^{-4}).               \tag{P-F52m}
\]
For every $q>0$, $HZ$ is an actual strictly positive central
operator on the same tensor space, since $H(q)>0$.
On the descended two-dimensional space its determinant polynomial is
\[
             \det(I+tHZ)=1+H(q)T(q)t+H(q)^2N(q)t^2,             \tag{P-F52n}
\]
where $T,N$ are precisely (P-F39c).
In the symmetric-vector notation $[v]$ of P-F39, direct convolution
gives the complete coefficient certificates
\[
 \begin{split}
 HT=[\alpha^H],\quad
 \alpha^H={}&(390,336,242,160,139,132,128,112,103,88,70,44,23,8,2),\\
 H^2N=[\beta^H],\quad
 \beta^H={}&(169556,164696,152852,139200,127656,118312,109695,\\
 &100480,90558,79776,68121,55832,44025,33832,25905,20000,\\
 &15558,12016,9093,6632,4587,2936,1699,864,377,136,39,8,1).
 \end{split}                                                    \tag{P-F52o}
\]
Every entry is positive, and each list is nonincreasing.
These are exact finite identities: if $h_j$ are the coefficients
in (P-F52m), extended symmetrically and by zero, then
\[
 \alpha_j^H=\sum_i h_i t_{j-i},\qquad
 \beta_j^H=\sum_{i,k}h_ih_kn_{j-i-k},                           \tag{P-F52p}
\]
using the explicitly displayed coefficients $t_j,n_j$ of (P-F39c).
The two coefficient sums at one are $81\cdot44=3564$ and
$81^2\cdot448=2939328$. The telescoping formula (P-F39h)
therefore realizes each as a nonnegative direct sum of the same
Tate-twisted projective-space Lefschetz characters.

Nevertheless $H$ is not divisible by $F$ in
$\mathbb C[q,q^{-1}]$. For a nonzero Laurent polynomial define
its width as its largest exponent minus its smallest exponent.
The extremal coefficients in a product are nonzero products,
so widths add. Here $F$ has width $32$ by (P-F39a), whereas
$H$ has width $8$, excluding $H=FJ$ for any nonzero Laurent
polynomial $J$. In particular these conditions do not even
force the prescribed factor to divide the normalization.

This proves the precise failure of characterization by the stated
finite operator and determinant properties. It does not assert
that $H$ has the source's additional lattice, integral-form, or
canonical-basis compatibility, or that it normalizes all other
coefficients. Those compatibility conditions remain the missing
datum; positivity of this actual action and its two Lefschetz
coefficients alone cannot supply them.\hfill$\square$

\paragraph{The exact spectral-denominator provenance of $F$.}
The additional rule that specifies this source normalization can
itself be identified explicitly. In GCT VII, p.~27,
equations~(32)--(35), the actual two-tensor braid operator
$\widehat R$ has eigenvalues
\[
       q_1=q^{9/2},\qquad q_2=-q^{-3/2},\qquad
       q_3=q^{-11/2},\qquad q_4=-q^{-15/2}.                    \tag{P-F52q}
\]
The positive spectral projector is expressed by Lagrange interpolation:
\[
  P_+=\sum_{i\in\{1,3\}}
       \frac{\prod_{j\ne i}(\widehat R-q_j)}
            {\prod_{j\ne i}(q_i-q_j)}.
\]
The source specifies the product of the two displayed denominators,
not the smallest common denominator after matrix cancellations:
\[
 f_p=\prod_{i\in\{1,3\}}\prod_{j\ne i}(q_i-q_j).
                                                               \tag{P-F52r}
\]
In particular the factors $q_1-q_3$ and $q_3-q_1$ both occur.
Their product is $-q^{-11}(q^{10}-1)^2$. The other four factors
are respectively
\[
 q^{-3/2}(q^6+1),\quad q^{-15/2}(q^{12}+1),\quad
 q^{-11/2}(q^4+1),\quad q^{-15/2}(q^2+1).
\]
Thus their exact product and its ratio to the convenient matrix
normalizer $f=q^{-5}(q^4+1)(q^6+1)$ are
\[
 \begin{split}
 f_p&=-q^{-33}(q^{10}-1)^2(q^6+1)(q^{12}+1)(q^4+1)(q^2+1),\\
 f_p/f&=-q^{-28}(q^{10}-1)^2(q^{12}+1)(q^2+1).
 \end{split}
\]
The source then removes the displayed vanishing factor by defining
$\widehat f_p=-q f_p/((q-1)^2f)$ on printed p.~28.
Consequently the centering used in P-F39 gives exactly
\[
 \begin{split}
 F=q^{11}\widehat f_p
  &=q^{-16}(1+q+\cdots+q^9)^2(1+q^2)(1+q^{12}).
 \end{split}                                                    \tag{P-F52s}
\]
This is a Laurent identity, not an evaluation at a singular
interpolation parameter. The factorization cancels $(q-1)^2$
and extends regularly to $q=1$, where $F(1)=10^2\cdot2\cdot2=400$.
It proves positivity of the coefficients directly. Expanding it
also gives exactly the complete coefficient list in (P-F39a).

Thus the finite provenance of $F$ is fixed once one supplies the
source's particular interpolation-denominator product and its
specified removal and centering factors. The positive central
action alone does not impose that interpolation rule, as the
alternative $H$ proves. This dictionary also does not assert that
the spectral-denominator rule supplies a categorical integral form,
canonical-basis multiplication, or the required geometric
intertwiner for all nonstandard coefficients. Those additional
compatibilities are distinct from the explicit algebraic
provenance just established.\hfill$\square$

\paragraph{An ordinary lambda-operation obstruction for the two coefficient models.}
The separate positive coefficient models cannot be combined into
the first two ordinary exterior-power characters of one graded
vector space, even after replacing $F$ by any common scalar
normalization. More precisely, there is no rational function
$J(q)$ regular and nonzero at $q=1$ and no finite-dimensional
$\mathbb Z$-graded vector space $V$ for which
\[
 \operatorname{ch}_q V=J(q)T(q),\qquad
 \operatorname{ch}_q\bigwedge^2V=J(q)^2N(q).
                                                        \tag{P-F52t}
\]
The same obstruction holds for a virtual graded-vector-space
class with the ordinary special lambda-ring operations.

Indeed evaluation of a graded character at one is its dimension.
Writing $r=44J(1)$ and using $T(1)=44$, $N(1)=448$ from
P-F39, the two equalities in (P-F52t) would give
\[
 r\in\mathbb Z,\quad r\ne0,\qquad
 \frac{r(r-1)}2=448J(1)^2=\frac{28}{121}r^2.
                                                        \tag{P-F52u}
\]
Cancelling $r$ yields $65r=121$, impossible for an integer.
For virtual classes the same calculation applies: their rank
is an integer and the rank of the second lambda operation is
$\binom r2$, as follows by expanding
$(1+t)^{r}$ (including negative integral exponents).
For the prescribed $F$ the numerical discrepancy is particularly
explicit:
\[
 \binom{17600}{2}=154871200\ne71680000=N_1(1).
                                                        \tag{P-F52v}
\]
Thus individual Tate--Lefschetz character realizations of
$T_1$ and $N_1$ do not give this common-object lambda
dictionary.

There is a related obstruction when $q$ is interpreted as
finite-field cardinality and the same Frobenius is required
over all finite extensions. Fix any integer $p>1$. There is
no single finite-dimensional linear operator $A$ over a
characteristic-zero field whose traces satisfy
\[
 \operatorname{Tr}(A^r)=T_1(p^r),\qquad
 \operatorname{Tr}\bigl(\bigwedge^2 A^r\bigr)=N_1(p^r)
                   \quad\hbox{for every }r\ge1.
                                                        \tag{P-F52w}
\]
To prove this, triangularize $A$ over an algebraic closure.
If its diagonal entries are $a_i$, summing
$a_i^r a_j^r$ over $i<j$ gives the elementary identity
\[
 2\operatorname{Tr}\bigl(\bigwedge^2 A^r\bigr)
       =\operatorname{Tr}(A^r)^2-\operatorname{Tr}(A^{2r}).
\]
The proposed traces would therefore make the Laurent
polynomial
\[
       2N_1(q)-T_1(q)^2+T_1(q^2)
                                                        \tag{P-F52x}
\]
vanish at every $q=p^r$. These are infinitely many distinct
nonzero values, so after multiplication by a power of $q$
the resulting polynomial must vanish identically. Evaluating
at one contradicts (P-F52v). No purity or semisimplicity
assumption is used in this contradiction.

These statements exclude the specified ordinary lambda
dictionary and the specified extension-compatible Frobenius
dictionary. A graded trace of an additional endomorphism,
a different relationship between the two coefficient objects,
or a different operation is not identified with an ordinary
graded dimension. In particular the actual central operator
and its determinant remain valid; the obstruction concerns
their proposed common-object cohomological interpretation,
not the algebraic determinant identity.\hfill$\square$
